\documentclass{article}
\usepackage[margin=28mm]{geometry}
\usepackage{amsmath,amssymb}
\usepackage{xurl}
\usepackage[hidelinks]{hyperref}
% Shared commands for notes in docs/paper-gaps/.

\DeclareUnicodeCharacter{2080}{\ifmmode{}_{0}\else\textsubscript{0}\fi}
\DeclareUnicodeCharacter{2081}{\ifmmode{}_{1}\else\textsubscript{1}\fi}
\DeclareUnicodeCharacter{2082}{\ifmmode{}_{2}\else\textsubscript{2}\fi}
\DeclareUnicodeCharacter{2099}{\ifmmode{}_{n}\else\textsubscript{n}\fi}

\newcommand{\Lean}{\textsc{Lean}}
\ExplSyntaxOn
\NewDocumentCommand{\leanid}{m}{
  \ifmmode
    \textsf{\detokenize{#1}}
  \else
    \group_begin:
    \urlstyle{sf}
    \tl_set:Nn \l_tmpa_tl {#1}
    \tl_replace_all:Nnn \l_tmpa_tl {\_} {_}
    \exp_args:NV \path \l_tmpa_tl
    \group_end:
  \fi
}
\ExplSyntaxOff
\providecommand{\tr}{\operatorname{tr}}
\newcommand{\tnleanIssueBase}{https://github.com/LionSR/TNLean/issues/}
\newcommand{\tnleanPrBase}{https://github.com/LionSR/TNLean/pull/}
\newcommand{\tnleanDocBase}{https://sirui-lu.com/TNLean/docs/}
\newcommand{\ghissue}[1]{\href{\tnleanIssueBase#1}{GitHub issue \##1}}
\newcommand{\ghpr}[1]{\href{\tnleanPrBase#1}{PR \##1}}
\newcommand{\doclink}[1]{\href{\tnleanDocBase#1.html}{\path{#1}}}

% Dirac notation and the identity operator, as in blueprint/src/macros/common.tex.
% \providecommand keeps note-local definitions authoritative.
\usepackage{dsfont}
\providecommand{\ket}[1]{|#1\rangle}
\providecommand{\bra}[1]{\langle #1|}
\providecommand{\braket}[2]{\langle #1 | #2 \rangle}
\providecommand{\ketbra}[2]{|#1\rangle\!\langle #2|}
\providecommand{\Id}{\mathds{1}}
\providecommand{\one}{\mathds{1}}
\providecommand{\C}{\mathbb{C}}
\providecommand{\MN}[1]{M_{#1}(\C)}
\providecommand{\MD}{M_D(\C)}

% Verdict marker: \gapnote{<kind>}{<status>}. Kind is the policy
% classification (clarification, local-correction, scope-restriction,
% unfaithful, false-source, open-gap); status is open, wip (elimination
% underway), resolved, or historical. Machine-read by the site index;
% prints a verdict line.
\newcommand{\gapnote}[2]{\par\noindent\textbf{Verdict:} #1 (#2).\par}

\providecommand{\gapnote}[2]{\par\noindent\textbf{Verdict:} #1 (#2).\par}
\title{The Supercell in Hastings' Translation-Invariant Case}
\author{}
\date{}
\begin{document}
\maketitle
\gapnote{clarification}{resolved}

\section{The source argument}
Hastings, \emph{Solving gapped Hamiltonians locally}~\cite{Hastings2006Solving},
\href{https://arxiv.org/abs/cond-mat/0508554v3}{arXiv:cond-mat/0508554v3},
source \path{proj3new.tex}, lines 1603--1642, considers a translationally
invariant system on a $d$-dimensional lattice with $n$ ground states, and the
local terms $M_i$ with their compressions $P M_i P$ to the ground space. Here
$M_i$ is the truncated filtered term of eq.~\texttt{mid} (lines 285--307),
within the error of eq.~\texttt{apb} of the filtered term $\tilde H_i$.
For $H=\sum_kh_k$ with gap $\Delta E$ above $E_0$ and a width $q>0$,
\begin{align}
    \tilde H_i=\frac{\Delta E}{\sqrt{2\pi q}}\int dt\,e^{iHt}h_ie^{-iHt}
    e^{-(t\Delta E)^2/(2q)},
    \label{eq:hastings05_translation_supercell_filtered_term}
\end{align}
and $M_i$ evolves $h_i$ under the sum $H_{loc}$ of
the $h_k$ near $i$ instead of $H$. Both are those of the note
\path{hastings05_filtered_term_approximation.tex} on the filtered terms, where
$\tilde H_i$ is written $\tilde H_i^0$. Choosing the ground states as eigenvectors of the translations, with
momenta $\vec k_a$, the matrix elements satisfy
$(M_{\vec x,y})_{ab}=(M_{\vec x+\vec x',y})_{ab}
    \exp[i\vec x'\cdot(\vec k_a-\vec k_b)]$ (lines 1613--1618).

The source then assumes the ``simplest case'' (lines 1620--1625): every
component of every momentum is a rational multiple $2\pi p/q$ whose
denominator $q$ does not depend on the system size. From this it derives a
supercell, whose size in direction $m$ is the least common multiple of the
denominators $q^m_a$, such that ``all of the $n$ ground states are unchanged by
translation by a supercell'' (lines 1626--1630). For a translation $\vec x'$
by whole supercells, $[M_{\vec x_1,y_1},M_{\vec x_2,y_2}]$ and
$[M_{\vec x_1,y_1},M_{\vec x_2+\vec x',y_2}]$ then agree in the low-energy
sector, and choosing $\vec x'$ to place the second site at a distance of order
the linear size of the system makes the commutator exponentially small in that
size (lines 1634--1641).

\section{The former restriction}
The first formalization did not derive the supercell from the momenta. It
assumed directly a subgroup $S$ of the symmetry group whose site permutations
$P_g$ leave every ground vector unchanged, $P_gv=v$ for $g\in S$ and $v$ in
the ground space, which is the conclusion of lines 1626--1630. From this it
proved, for covariant terms $H_{g\cdot k}=P_gH_kP_g^\dagger$ and an invariant
Hamiltonian,
\begin{align}
    PM_{g\cdot j}P & =PM_jP,
    \label{eq:hastings05_translation_supercell_compression} \\
    [PM_iP,PM_jP]  & =[PM_iP,PM_{g\cdot j}P]\qquad(g\in S),
    \label{eq:hastings05_translation_supercell_commutator}
\end{align}
and, in the setting of eq.~\texttt{blecm}, the decay of the low-energy
commutator at any separation $l$ of the supports of $H_i$ and
$H_{g\cdot j}$.\footnote{\leanid{TNLean.Lattice.groundCompression_sitePermMatrix_conj},
    \leanid{TNLean.Lattice.groundCompression_commutator_eq_of_mem_supercell},
    \leanid{TNLean.Lattice.norm_groundCompression_commutator_le_of_mem_supercell}
    in \doclink{TNLean/Lattice/TranslationSupercell}, and
    \leanid{TNLean.Lattice.norm_groundCompression_commutator_le_of_supercell_separation}
    in \doclink{TNLean/Lattice/TranslationSupercellDecay}.}
These results still take $S$ as a hypothesis; they are the steps of the
argument below.

\section{Removing the supercell assumption}
The restriction was removed in two steps.

\emph{The supercell from the momenta.} For an abelian group $G$, such as the
translations of a lattice, the operators $P_g$ commute with an invariant $H$
and with one another, and each has finite order, so they are simultaneously
diagonalizable on the ground space. The ground space is the sum of its
momentum sectors $\{v:P_gv=\chi(g)v\text{ for all }g\}$ over characters
$\chi$, and the intersection $S$ of the kernels of the characters with a
nonzero sector is exactly the subgroup fixing every ground
vector.\footnote{\leanid{TNLean.Lattice.iSup_momentumSector},
    \leanid{TNLean.Lattice.mem_supercell_iff} and
    \leanid{TNLean.Lattice.sitePermMatrix_mulVec_eq_of_mem_supercell} in
    \doclink{TNLean/Lattice/TranslationMomentum}.}

\emph{The spacing of the supercell.} The source's hypothesis of lines
1620--1625 says that each component of each momentum is $2\pi p/q$ with a
denominator independent of the system size; with $N$ the least common
multiple of the denominators, $\chi(g)^N=1$ for every character with a nonzero
sector, so every $N$-th power lies in $S$. For a lattice with $n$ cells along
one direction, let $\pi$ be the cell coordinate in that direction, with every
edge of the lattice graph joining cells at most a positive integer $r$ apart,
and let $e$ be the unit translation in that direction. The circular distance
$\min(a,n-a)$ of the coordinate difference $a$ is at most $r$ times the graph
distance $d$, and for any sites $x,y$ some translation $g=(e^t)^N$ satisfies
$\lfloor n/2\rfloor<r\,d(x,g\cdot y)+N$. For a ring, and in the composite
bound below, $r=1$, so $g$ carries any site to graph distance more than
$\lfloor n/2\rfloor-N$ from any
other.\footnote{\leanid{TNLean.Lattice.pow_mem_supercell},
    \leanid{SimpleGraph.natAbs_valMinAbs_sub_le_edist},
    \leanid{TNLean.Lattice.exists_mem_supercell_lt_edist_add} and its ring case
    \leanid{TNLean.Lattice.exists_mem_supercell_lt_edist_add_ring} in
    \doclink{TNLean/Lattice/TranslationSupercellSpacing}.}

\emph{The composite bound.} For terms indexed by the sites, with supports of
diameter at most $R$ containing their sites, covariant under an abelian
translation group with momenta of denominator $N$, and a gap above $E_0$, the
low-energy commutator at $i$ and $j$ obeys the bound of eq.~\texttt{blecm} at
every separation $l$ with $2R+l+N\le\lfloor n/2\rfloor$. Taking
$l=\lfloor n/2\rfloor-N-2R$ makes it exponentially small in $n$. This is first
proved for the filtered terms $\tilde H_i$, whose compressions equal those of
the $h_i$. For the source's truncated terms $M_i$ of eq.~\texttt{mid}, within
$\varepsilon_i$ of $\tilde H_i$ by eq.~\texttt{apb}, the commutator moves by at
most $2\varepsilon_i\lVert h_j\rVert+2\lVert h_i\rVert\varepsilon_j$.\footnote{\leanid{TNLean.Lattice.norm_groundCompression_commutator_le_of_momenta},
    \leanid{TNLean.Lattice.norm_groundCompression_commutator_le_of_norm_sub_le} and
    \leanid{TNLean.Lattice.norm_groundCompression_commutator_truncated_le_of_momenta}
    in \doclink{TNLean/Lattice/TranslationSupercellCommutator}.}
No hypothesis beyond the source's is assumed.

\section{The truncated terms}
The errors $\varepsilon_i$ of eq.~\texttt{apb} depend on the truncation radius
$l_{proj}$ and not on $n$, so the bound above for the truncated $M_i$ is not
exponentially small in $n$. The source avoids this by applying the Gaussian
filter a second time to $M_i$ itself: the filter $\tilde M_i^0$ of $M_i$
under $H$ satisfies
\begin{align}
    P\tilde M_i^0P=PM_iP,
    \label{eq:hastings05_translation_supercell_second_filter}
\end{align}
and the
low-energy commutator of the $M_i$ obeys eq.~\texttt{blecm} directly
(lines 1555--1601). This is formalized. Each $M_i$ acts on a set $K_i\ni i$ of
diameter at most $R'$; with covariant truncations, $s_{g\cdot k}=g\cdot s_k$,
one has $M_{g\cdot k}=P_gM_kP_g^\dagger$, so $PM_{g\cdot j}P=PM_jP$ for $g$ in
the supercell. Eq.~\texttt{blecm} applied to $M_i$ and $M_{g\cdot j}$ then
bounds the low-energy commutator of the $M_i$ at every separation $l$ with
$2R'+l+N\le\lfloor n/2\rfloor$, with no truncation error, which gives the
conclusion of lines 1634--1641 for the source's
$M_i$.\footnote{\leanid{TNLean.Lattice.truncatedTerm_smul} and
    \leanid{TNLean.Lattice.norm_groundCompression_truncatedTerm_commutator_le_of_momenta}
    in \doclink{TNLean/Lattice/TranslationSupercellTruncated}.}

\section{A slip}
At line 1638 the source writes the distance ``between $(\vec x_1,y_2)$ and
$(\vec x_2+\vec x',y_2)$''; the first site is $(\vec x_1,y_1)$.

\section{Arbitrary momentum denominators}

Lines 1644--1654 assert exponential smallness even when the momentum
denominators grow with the system size, but omit the proof. A
finite-dimensional argument proves this assertion without a condition on
the momenta.

Let $K$ be the actual ground eigenspace, let $P$ be its orthogonal
projection, and put $C_K(W)=(PWP)|_K$. A physical translation preserving
$H$ restricts to a unitary $u$ on $K$. It commutes with $P$, and hence
\[
 T^k(C_K(W))=C_K(U^kW(U^k)^*),\qquad T(B)=uBu^{-1}.
\]
Here $T$ preserves the operator norm on the $n^2$-dimensional space of
operators on $K$, where $n=\dim K$. The ambient compressed commutator
has exactly the norm of its restriction to $K$.\footnote{Declarations:
\leanid{TNLean.Lattice.groundTranslation},
\leanid{TNLean.Lattice.groundTranslationConjugation},
\leanid{TNLean.Lattice.groundTranslationConjugation_pow_compression}, and
\leanid{TNLean.Lattice.norm_groundSpaceCompression_commutator_eq}.}

More generally, let $T$ preserve a norm on a complex space of dimension
$d$. Its characteristic polynomial $p(z)=\sum_{k=0}^d p_kz^k$ is monic,
its roots have modulus one, and consequently
\[
 |p_0|=1,\qquad \sum_{k=1}^d|p_k|\le2^d-1.
\]
Cayley--Hamilton gives $1=-\sum_{k=1}^d p_0^{-1}p_kT^k$. Thus
\[
 \|[A,B]\|\le(2^{n^2}-1)
       \max_{1\le k\le n^2}\|[A,T^k(B)]\|.
\]
The characteristic-polynomial estimate uses no separation of the
translation eigenvalues.\footnote{Declaration:
\leanid{TNLean.Lattice.norm_commutator_le_of_finite_translation_orbit_bounds}.}

Choose an integer $m\ge\max\{1,n^2\}$ independent of the system size,
let $N$ be the number of cells in one direction, and put
$L=\lfloor N/[2(m+1)]\rfloor$. If two cells have circular separation at
most $\lfloor L/2\rfloor$, translating the second by $kL$ cells, for
$1\le k\le m$, gives separation at least $\lfloor L/2\rfloor$.
Indeed, a signed representative of the initial displacement has absolute
value at most $\lfloor L/2\rfloor$, and adding $kL$ puts it between
$\lfloor L/2\rfloor$ and $N-\lfloor L/2\rfloor$. If the initial cells
are already farther apart, no translation is needed.\footnote{Declarations:
\leanid{TNLean.Lattice.natAbs_valMinAbs_add_mul_ge_of_near} and
\leanid{TNLean.Lattice.edist_pow_mul_ge_of_near}.}

For the local observables, suppose the graph radius is at most $\rho$,
the support cardinality is at most $v_{\rm obs}$, and their norms are at
most $J$. The propagation and original-gap estimate gives a positive
spatial rate $\gamma$ and the factor
$C_0=(4v_{\rm obs}+4)J^2$. It applies to arbitrary observables on their
supports, rather than only to members of the chosen local decomposition.
Apply it to the above translates of the second observable. The finite
orbit estimate, the fixed support offset, and
\[
 \frac{N}{4(m+1)}\le\frac{L}{2}+\frac12
\]
give the uniform bound
\[
 \|[PWP,PBP]\|\le
 (2^m-1)C_0 e^{\gamma(2\rho+2)}
       e^{-\gamma N/[4(m+1)]}.
\]
For small $N$, including $L=0$, the same bound follows from
$\|[PWP,PBP]\|\le2J^2$; its displayed prefactor absorbs these cases.
The integer rounding and the small-size comparison are both
proved.\footnote{Declarations:
\leanid{TNLean.Lattice.translation_half_natDiv_lower_bound},
\leanid{TNLean.Lattice.translation_exp_floor_decay_le}, and
\leanid{TNLean.Lattice.translation_finiteOrbit_decay_of_near_far}.}

Only the Hamiltonian must be invariant. The physical translations
preserve the graph geometry, so conjugation preserves support radius,
cardinality, and norm. The local terms of an arbitrary decomposition
need not be covariant, and neither do their truncation sets: the proof
translates the second observable itself. In particular it applies to
the source's actual $M_i$, whose support and norm were proved from the
truncated Gaussian construction. The constants are uniform when the
local data and the ground-space dimension are uniformly bounded, as
assumed in the source passage.
The physical bounds and their specialization to the
actual truncated terms are formalized without a supplied spatial-decay
premise.\footnote{Declarations:
\leanid{TNLean.Lattice.norm_groundCompression_commutator_le_of_finite_observable_orbit},
\leanid{TNLean.Lattice.norm_groundCompression_commutator_le_finite_translation_exp}, and
\leanid{TNLean.Lattice.norm_groundCompression_truncated_commutator_le_finite_translation_exp},
in \doclink{TNLean/Lattice/TranslationFiniteOrbitCommutator}.}

\bibliographystyle{alpha}
\bibliography{references}

\end{document}
